% ===== CHƯƠNG 2: CƠ SỞ LÝ THUYẾT =====

\chapter{Cơ Sở Lý Thuyết}

\section{Bài toán tìm đường trên đồ thị có trọng số}

\subsection{Định nghĩa chính thức}

Cho một đồ thị vô hướng có trọng số $G = (V, E, w)$ với:
\begin{itemize}
  \item $V$ là tập hợp các đỉnh (vertices) --- trong bối cảnh này là các ga metro.
  \item $E$ là tập hợp các cạnh (edges) --- đoạn đường nối hai ga liền kề.
  \item $w: E \to \mathbb{R}^+$ là hàm trọng số gán cho mỗi cạnh một giá trị dương (chi phí).
\end{itemize}

Bài toán tìm đường ngắn nhất từ ga nguồn $s$ đến ga đích $g$ là tìm một đường đi $P = [s = v_0, v_1, \ldots, v_n = g]$ sao cho tổng chi phí:
\[
  \text{cost}(P) = \sum_{i=0}^{n-1} w(v_i, v_{i+1})
\]
là nhỏ nhất trong tất cả các đường đi từ $s$ đến $g$.

\subsection{Không gian trạng thái (State Space)}

Mỗi trạng thái trong quá trình tìm kiếm được định nghĩa là:
\[
  \text{State} = (\text{station\_name}, \text{cumulative\_cost}, \text{path}, \text{explored\_count})
\]

Trạng thái bắt đầu: $s_0 = (s, 0, [s], 0)$

Trạng thái mục tiêu: Bất kỳ trạng thái nào có station\_name = $g$

\subsection{Hàm hành động (Actions)}

Từ một trạng thái hiện tại, có thể chuyển tới các trạng thái lân cận bằng cách đi qua các cạnh từ ga hiện tại tới các ga hàng xóm.

\section{Biểu diễn đồ thị trong ứng dụng}

Ứng dụng sử dụng biểu diễn **danh sách kề** (adjacency list). Cấu trúc dữ liệu chính:

\begin{lstlisting}[language=Python, caption=Cấu trúc dữ liệu Station và Edge (metro\_app/data.py)]
@dataclass(frozen=True)
class Station:
    name: str
    lat: float
    lon: float

@dataclass(frozen=True)
class Edge:
    target: str
    travel_time: float
    line: str
    distance_km: float
\end{lstlisting}

Đồ thị được lưu dưới dạng:
\[
  \text{graph: dict[str, list[Edge]]}
\]

Với mỗi ga (key), lưu danh sách các cạnh tới các ga hàng xóm.

\section{Uniform Cost Search (UCS)}

\subsection{Ý tưởng chính}

UCS là một thuật toán tìm kiếm hướng mục tiêu (goal-directed) sử dụng hàng đợi ưu tiên (priority queue). Thuật toán luôn chọn trạng thái có chi phí nhỏ nhất để khám phá, đảm bảo rằng lần đầu tiên đạt được mục tiêu là lần với chi phí nhỏ nhất.

\subsection{Pseudocode}

\begin{lstlisting}[caption=Uniform Cost Search Pseudocode]
function UniformCostSearch(graph, start, goal):
  queue = PriorityQueue()
  queue.push((0, start))

  came_from = {start: None}
  best_cost = {start: 0}
  explored = 0

  while queue is not empty:
    current_cost, current = queue.pop()

    // Skip if already found better path
    if current_cost > best_cost[current]:
      continue

    explored += 1

    // Found goal
    if current == goal:
      path = reconstruct_path(came_from, goal)
      return (path, current_cost, explored)

    // Expand neighbors
    for edge in graph[current]:
      new_cost = current_cost + edge.travel_time
      if edge.target not in best_cost or new_cost < best_cost[edge.target]:
        best_cost[edge.target] = new_cost
        came_from[edge.target] = current
        queue.push((new_cost, edge.target))

  return FAILURE
\end{lstlisting}

\subsection{Tính chất}

\begin{itemize}
  \item \textbf{Hoàn chỉnh (Complete):} Có. Nếu tồn tại đường đi, UCS sẽ tìm thấy.
  \item \textbf{Tối ưu (Optimal):} Có. Lần đầu pop goal từ queue là lần với chi phí nhỏ nhất.
  \item \textbf{Heuristic:} Không dùng heuristic.
  \item \textbf{Độ phức tạp thời gian:} $O(b^{C^*/\epsilon})$ với $b$ là branching factor, $C^*$ là chi phí tối ưu, $\epsilon$ là chi phí cạnh nhỏ nhất.
  \item \textbf{Độ phức tạp không gian:} $O(b^{C^*/\epsilon})$ --- lưu tất cả node trong queue.
\end{itemize}

\section{Thuật toán Dijkstra}

\subsection{Ý tưởng chính}

Dijkstra là một trường hợp đặc biệt của UCS khi heuristic $h = 0$. Tuy nhiên, trong lý thuyết truyền thống, Dijkstra được triển khai khác biệt với mục tiêu tính toàn bộ **Single-Source Shortest Path (SSSP)** từ một nguồn đến tất cả các đỉnh khác, không dừng sớm khi tìm được một mục tiêu cụ thể.

\subsection{Pseudocode}

\begin{lstlisting}[caption=Dijkstra Algorithm Pseudocode]
function Dijkstra(graph, start, goal):
  distances = {node: infinity for all nodes}
  distances[start] = 0

  came_from = {start: None}
  unvisited = {all nodes}
  explored = 0

  while unvisited is not empty:
    current = node in unvisited with min distances[node]

    if distances[current] == infinity:
      break  // No path exists

    unvisited.remove(current)
    explored += 1

    // Continue to compute full SSSP (no early exit)
    for edge in graph[current]:
      neighbor = edge.target
      new_distance = distances[current] + edge.travel_time

      if new_distance < distances[neighbor]:
        distances[neighbor] = new_distance
        came_from[neighbor] = current

  // Reconstruct path to goal using came_from
  path = reconstruct_path(came_from, goal)
  return (path, distances[goal], explored)
\end{lstlisting}

\subsection{Tính chất}

\begin{itemize}
  \item \textbf{Hoàn chỉnh:} Có.
  \item \textbf{Tối ưu:} Có.
  \item \textbf{Heuristic:} Không.
  \item \textbf{Độ phức tạp thời gian:} $O((V + E) \log V)$ với priority queue.
  \item \textbf{Độ phức tạp không gian:} $O(V)$ --- lưu distances table.
\end{itemize}

\subsection{UCS vs Dijkstra}

UCS và Dijkstra đều tính toàn bộ chi phí tối ưu, nhưng có sự khác biệt:

\begin{itemize}
  \item \textbf{UCS:} Goal-directed. Dừng ngay khi pop mục tiêu từ queue. Số node explored tối thiểu.
  \item \textbf{Dijkstra:} SSSP-oriented. Tính toàn bộ distances từ một nguồn đến tất cả node. Số node explored = toàn bộ graph reachable từ start, thường > UCS.
\end{itemize}

Trong báo cáo này, Dijkstra được cài đặt với early-exit (dừng khi pop goal), nhưng đạt được SSSP trước khi reconstruct path, nên số node explored cao hơn UCS.

\section{Greedy Best-First Search}

\subsection{Ý tưởng chính}

Greedy chỉ dùng heuristic estimate (không dùng chi phí thực), luôn chọn node có $h(n)$ nhỏ nhất (gần mục tiêu nhất theo ước lượng). Thuật toán nhanh nhưng không đảm bảo tối ưu.

\subsection{Pseudocode}

\begin{lstlisting}[caption=Greedy Best-First Search Pseudocode]
function GreedyBestFirstSearch(graph, stations, start, goal):
  queue = PriorityQueue()
  h_start = heuristic(stations, start, goal)
  queue.push((h_start, 0, start))

  came_from = {start: None}
  visited = {set of visited nodes}
  explored = 0

  while queue is not empty:
    _, current_cost, current = queue.pop()

    if current in visited:
      continue

    visited.add(current)
    explored += 1

    if current == goal:
      path = reconstruct_path(came_from, goal)
      return (path, current_cost, explored)

    for edge in graph[current]:
      neighbor = edge.target
      if neighbor not in visited:
        h_neighbor = heuristic(stations, neighbor, goal)
        new_cost = current_cost + edge.travel_time
        came_from[neighbor] = current
        queue.push((h_neighbor, new_cost, neighbor))

  return FAILURE
\end{lstlisting}

\subsection{Tính chất}

\begin{itemize}
  \item \textbf{Hoàn chỉnh:} Có (với graph hữu hạn và tracking visited).
  \item \textbf{Tối ưu:} Không. Greedy có thể tìm thấy đường không phải chi phí nhỏ nhất.
  \item \textbf{Heuristic:} Có.
  \item \textbf{Độ phức tạp thời gian:} $O(b^m)$ với $m$ là độ sâu bài toán.
  \item \textbf{Độ phức tạp không gian:} $O(b^m)$.
\end{itemize}

\section{Thuật toán A*}

\subsection{Ý tưởng chính}

A* kết hợp chi phí thực $g(n)$ (từ start tới $n$) và ước lượng $h(n)$ (từ $n$ tới goal) qua công thức:
\[
  f(n) = g(n) + h(n)
\]

A* luôn khám phá node có $f(n)$ nhỏ nhất, đảm bảo tối ưu nếu $h(n)$ là admissible (không overestimate).

\subsection{Pseudocode}

\begin{lstlisting}[caption=A* Algorithm Pseudocode]
function AStar(graph, stations, start, goal):
  queue = PriorityQueue()
  g_start = 0
  h_start = heuristic(stations, start, goal)
  queue.push((g_start + h_start, g_start, start))

  came_from = {start: None}
  best_g = {start: 0}
  explored = 0

  while queue is not empty:
    _, g_current, current = queue.pop()

    if g_current > best_g[current]:
      continue

    explored += 1

    if current == goal:
      path = reconstruct_path(came_from, goal)
      return (path, g_current, explored)

    for edge in graph[current]:
      neighbor = edge.target
      g_neighbor = g_current + edge.travel_time

      if neighbor not in best_g or g_neighbor < best_g[neighbor]:
        best_g[neighbor] = g_neighbor
        came_from[neighbor] = current
        h_neighbor = heuristic(stations, neighbor, goal)
        f_neighbor = g_neighbor + h_neighbor
        queue.push((f_neighbor, g_neighbor, neighbor))

  return FAILURE
\end{lstlisting}

\subsection{Tính chất}

\begin{itemize}
  \item \textbf{Hoàn chỉnh:} Có.
  \item \textbf{Tối ưu:} Có (khi $h(n)$ admissible).
  \item \textbf{Heuristic:} Có.
  \item \textbf{Độ phức tạp thời gian:} $O(b^{C^*/\epsilon})$ --- có thể tốt hơn UCS nếu heuristic tốt.
  \item \textbf{Độ phức tạp không gian:} $O(b^{C^*/\epsilon})$.
\end{itemize}

\section{Hàm Heuristic và Tính Admissibility}

\subsection{Định nghĩa Admissibility}

Một heuristic $h(n)$ được gọi là **admissible** nếu nó không bao giờ overestimate chi phí thực từ $n$ tới goal:
\[
  h(n) \leq h^*(n) \quad \forall n
\]

với $h^*(n)$ là chi phí thực tế nhỏ nhất từ $n$ tới goal.

Nếu $h(n)$ admissible, A* được chứng minh là optimal.

\subsection{Hàm Heuristic trong ứng dụng}

Ứng dụng sử dụng:

\begin{lstlisting}[language=Python, caption=Hàm Heuristic (metro\_app/algorithms.py)]
def heuristic(stations: dict[str, Station], current: str, goal: str) -> float:
    here = stations[current]
    destination = stations[goal]
    return haversine_km(here.lat, here.lon, destination.lat, destination.lon) / 35.0 * 60.0
\end{lstlisting}

Công thức: $h(n) = \frac{\text{haversine}(n, goal)}{35} \times 60$

\subsection{Tính Admissibility}

Hàm này admissible vì:
\begin{enumerate}
  \item Haversine tính khoảng cách thẳng (đường tròn lớn) giữa hai điểm.
  \item Trong metro, đường đi thực tế luôn $\geq$ đường thẳng (do phải đi theo tuyến).
  \item Do đó: haversine $\leq$ khoảng cách thực tế.
  \item Chia cho 35 km/h (tốc độ metro) và nhân 60 (chuyển sang phút) cho ra giá trị time estimate luôn $\leq$ chi phí thực tế.
  \item Suy ra $h(n)$ admissible.
\end{enumerate}

\section{So sánh Lý Thuyết Bốn Thuật toán}

\begin{table}[H]
  \centering
  \caption{So sánh lý thuyết bốn thuật toán tìm đường}
  \label{tab:comparison}
  \begin{tabular}{lcccccc}
    \toprule
    \textbf{Tiêu chí} & \textbf{UCS} & \textbf{Greedy} & \textbf{A*} & \textbf{Dijkstra} \\
    \midrule
    Hoàn chỉnh (Complete) & ✓ & ✓* & ✓ & ✓ \\
    Tối ưu (Optimal) & ✓ & ✗ & ✓** & ✓ \\
    Dùng heuristic & ✗ & ✓ & ✓ & ✗ \\
    Time complexity & $O(b^{C^*/\epsilon})$ & $O(b^m)$ & $O(b^{C^*/\epsilon})$ & $O((V+E)\log V)$ \\
    Space complexity & $O(b^{C^*/\epsilon})$ & $O(b^m)$ & $O(b^{C^*/\epsilon})$ & $O(V)$ \\
    \bottomrule
  \end{tabular}
\end{table}

*Greedy hoàn chỉnh với graph hữu hạn và tracking visited.

**A* tối ưu khi $h(n)$ admissible.

\section{Kết luận chương}

Chương này đã trình bày chi tiết bốn thuật toán tìm đường, mỗi thuật toán có những ưu điểm và nhược điểm riêng. UCS và Dijkstra đảm bảo tối ưu nhưng không dùng heuristic. Greedy nhanh nhưng không tối ưu. A* cân bằng tốc độ và tính tối ưu nhờ vào heuristic admissible. Chương tiếp theo sẽ trình bày cách áp dụng những thuật toán này vào ứng dụng thực tế.
